% Chapter 13 of Optimizers in the Wild: weight decay.
\newcommand{\decayball}{\mathcal{B}_\lambda}
\newcommand{\decayeff}{\eta_{\mathrm{eff}}}
\newcommand{\decayfont}{\csname footnotesize\endcsname}

\casechapter{Weight Decay}{Frank--Wolfe, rotational equilibrium and the effective learning rate}{AdamW's decoupled weight decay}{ch:decay}
\begin{refsection}[chapters/13-decay.bib]

\begin{chaptersummary}
Pretraining recipes multiply the weights by $1-\eta\lambda$ before every AdamW step, with $\lambda=0.1$, although a weight matrix followed by a normalization layer computes the same function at every scale. When the step has a fixed length in a norm, decoupled decay is exactly the Frank--Wolfe method on the ball of radius $1/\lambda$: we prove that the ball attracts every initialization at rate $(1-\eta\lambda)^t$, that the learning-rate schedule acts as a Frank--Wolfe step-size schedule, and where in the ball consistent and noisy coordinates settle. When the loss is scale invariant, we prove that the norm converges from every initialization to an equilibrium at which the weights rotate by $\arccos(1-\eta\lambda)\approx\sqrt{2\eta\lambda}$ per step whatever the size of the gradients, so that only the product $\eta\lambda$ matters, and that the rotation lags any schedule. Momentum, preconditioning, gradients that respond to the effective learning rate, and the effects of decay on generalization are what these idealizations miss.
\end{chaptersummary}

\section{In the wild: a regularizer that does not regularize}\label{sec:decay:wild}

For stochastic gradient descent, weight decay and an $L_2$ penalty are the same thing: adding $\frac\lambda2\norm{\theta}_2^2$ to the loss adds $\lambda\theta$ to the gradient, and $\theta\leftarrow\theta-\eta(g+\lambda\theta)$ is $\theta\leftarrow(1-\eta\lambda)\theta-\eta g$. For Adam they differ, because the gradient of the penalty is divided by $\sqrt{\hat v}$ with the rest of the gradient. \Textcite{LoshchilovHutter2019} applied the shrinkage outside the adaptive step and found that this \term{decoupled weight decay}\index{weight decay!decoupled} makes the best decay less dependent on the learning rate and improves Adam's generalization; the result is AdamW. Every optimizer of this chapter has the form
\begin{equation}\label{eq:decay:update}
  \theta_{t+1}=(1-\eta_t\lambda)\,\theta_t-\eta_t\,u_t ,
\end{equation}
where $u_t$ is the optimizer's own step direction: the gradient for SGD, $\hat m_t/(\sqrt{\hat v_t}+\epsilon)$ for AdamW, the orthogonalized momentum for Muon. We say learning rate for $\eta_t$ in this section and step size in the theory; the two are the same.

\begin{practice}[GPT-3, \textcite{BrownEtAl2020}; LLaMA, \textcite{TouvronEtAl2023}; Muon at scale, \textcite{LiuEtAl2025Muon}]
Train with AdamW and decoupled weight decay $\lambda=0.1$. The peak learning rates of the GPT-3 models range from $6\times10^{-5}$ (the largest) to $6\times10^{-4}$ (the smallest), so each step shrinks the weights by a factor between $1-6\times10^{-6}$ and $1-6\times10^{-5}$; the factor follows the schedule, while $\lambda$ stays fixed. Decay the weight matrices; by a convention of open training code (folklore), biases and the gains of normalization layers are not decayed. When Muon replaces AdamW on the hidden matrices, keep a decoupled decay of the same form. PyTorch's AdamW multiplies the weights by $1-\eta_t\lambda$, with the default $\lambda=10^{-2}$.
\end{practice}

Two features of this practice do not fit the picture of a penalty. AdamW at $\lambda=0.1$ does not minimize the loss plus $\frac\lambda2\norm{\theta}_2^2$, which for a model with billions of weights would dwarf the loss; we shall see that it keeps the weights in the ball $\norm{\theta}_\infty\le1/\lambda=10$, far larger than the weights of a trained network. And almost every weight matrix of a transformer feeds a normalization layer, which erases the scale of its input.

\begin{observation}[\textcite{VanLaarhoven2017}; \textcite{ZhangEtAl2019}]\label{obs:decay:elr}
If the output of a weight matrix $W$ is normalized (by batch, layer or RMS normalization), the network computes the same function with $W$ and with $cW$ for every $c>0$, so an $L_2$ penalty on $W$ does not constrain the function. It changes the scale of $W$, and with it the size of each step relative to $W$: an effective learning rate. \Textcite{VanLaarhoven2017} made this argument for SGD and found that Adam removes the dependence on the scale only partly; \textcite{ZhangEtAl2019} identified the effective learning rate as one of three mechanisms by which weight decay improves the test error of vision networks.
\end{observation}

\begin{observation}[\textcite{WanEtAl2021}; \textcite{KossonMessmerJaggi2024}]\label{obs:decay:rotation}
With normalization layers and weight decay, the norm of each scale-invariant weight and its angle of rotation per step settle to a steady state. \Textcite{WanEtAl2021} derived the steady-state angle $\sqrt{2\eta\lambda}$ for SGD and confirmed their predictions on ImageNet and MS COCO. \Textcite{KossonMessmerJaggi2024} called the steady state \term{rotational equilibrium}\index{rotational equilibrium}, analysed and measured it for single neurons under AdamW, Lion and SGD with momentum, and argued that its balance across layers and neurons is part of why AdamW works better than Adam with an $L_2$ penalty.
\end{observation}

\begin{observation}[\textcite{LoshchilovHutter2019}; \textcite{WangAitchison2024}]\label{obs:decay:horizon}
The more batch passes, the smaller the optimal $\lambda$. \Textcite{LoshchilovHutter2019} proposed $\lambda=\lambda_{\mathrm{norm}}\sqrt{b/(BT)}$ for batch size $b$, $B$ training points and $T$ epochs, ``merely one possibility informed by few experiments''. \Textcite{WangAitchison2024} read AdamW's weights as an exponential moving average of its updates with a timescale of $1/(\eta\lambda)$ iterations, and report that the optimal timescale measured in epochs is roughly the same across model and dataset sizes.
\end{observation}

Muon taught the same lesson: \textcite{LiuEtAl2025Muon} name the addition of weight decay as one of the two changes that made Muon work at scale. A counter-observation: \textcite{LobachevaEtAl2021} found that with batch normalization and weight decay, training can be periodic rather than stationary; the norm shrinks, the effective learning rate grows, training destabilizes without diverging, and a new period begins.\sidenote{\Textcite{DAngeloEtAl2024} argue that for language models trained for about one epoch, weight decay acts on optimization rather than on generalization: it lowers the training loss and stabilizes training in low precision. The mechanisms of this chapter are optimization mechanisms.}

What, then, does $\lambda$ set in a training run? This chapter gives two answers, one for each of two idealizations: the radius of a ball that the weights cannot leave, and the speed at which scale-invariant weights rotate.

\section{The idealization: a shrinking step and a loss that ignores scale}\label{sec:decay:idealization}

The first idealization keeps the geometry of the step and nothing else.

\begin{assumption}[Steepest descent with decoupled decay]\label{ass:decay:nsd}
Fix a norm $\norm{\cdot}$ on $\R^d$ with dual norm $\norm{\cdot}_*$, a decay $\lambda>0$, and step sizes $\eta_t>0$ with $\eta_t\lambda\le1$. The iteration is
\[
  \theta_{t+1}=(1-\eta_t\lambda)\,\theta_t+\eta_t\lmo(g_t),
\]
that is, \eqref{eq:decay:update} with $u_t=-\lmo(g_t)$, where $g_t$ is whatever vector the optimizer feeds to its linear minimization oracle: a gradient, a stochastic gradient or a moving average of gradients.
\end{assumption}

AdamW without memory ($\beta_1=\beta_2=0$, $\epsilon=0$) is the $\ell_\infty$ case, $\lmo(g)=-\sign(g)$ (\cref{ch:sign}); Lion with decay is the $\ell_\infty$ case applied to an interpolation of momentum and gradient; Muon with decay is the spectral-norm case, $\lmo(G)=-\polar(G)$, applied to the momentum (\cref{ch:spectral}). Norms, oracles and dual norms are those of \cref{ch:steepest}. The second idealization keeps only the symmetry that normalization layers create.

\begin{definition}[Scale invariance]\label[definition]{def:decay:scaleinv}
A function $\Loss\colon\R^d\setminus\{0\}\to\R$ is \term{scale invariant}\index{scale invariance} if $\Loss(c\theta)=\Loss(\theta)$ for every $\theta\ne0$ and $c>0$.
\end{definition}

A layer normalization without gain satisfies $\mathrm{LN}(cWx)=\mathrm{LN}(Wx)$ for $c>0$, up to the small constant inside its square root, so a loss that depends on $W$ only through $\mathrm{LN}(Wx)$ is scale invariant in $W$; batch and RMS normalization, weight normalization and the normalization of queries and keys give the same structure. Throughout, $\norm{\theta}_2$ is the Euclidean norm of all entries of the weight in question (the Frobenius norm of a matrix) and $\hat\theta=\theta/\norm{\theta}_2$ its direction.

\begin{assumption}[Rotation on a sphere]\label{ass:decay:sphere}
The loss $\Loss$ and each per-example loss $\ell(\cdot;\xi)$ are scale invariant and differentiable; the step direction $u_t$ in \eqref{eq:decay:update} is orthogonal to $\theta_t$; and, in the main theorem, the stochastic gradient has a constant length on the unit sphere, $\norm{g_t}_2\norm{\theta_t}_2=\Gamma$.
\end{assumption}

\begin{dictionary}
the ball $\decayball=\{\theta:\norm{\theta}\le1/\lambda\}$ & the set AdamW's weights are drawn into: $\abs{\theta_i}\le10$ at $\lambda=0.1$\\
Frank--Wolfe step size $\gamma_t=\eta_t\lambda$ & the learning-rate schedule times the decay\\
the vertex $\lmo(g_t)/\lambda$ & where a step points: $-\sign(g_t)/\lambda$ for AdamW without memory\\
scale invariance, $\Loss(c\theta)=\Loss(\theta)$ & a weight matrix followed by a normalization layer\\
the direction $\hat\theta$; the norm $\norm{\theta}_2$ & what the layer computes; a free variable\\
$\Gamma=\norm{g}_2\norm{\theta}_2$ & the gradient norm of the same layer rescaled to unit norm\\
the angle $\omega_t$ between $\theta_t$ and $\theta_{t+1}$ & the rotation of the weight matrix in one step\\
$\arccos(1-\eta\lambda)\approx\sqrt{2\eta\lambda}$ & equilibrium rotation: $0.0077$ rad per step at $\eta=3\times10^{-4}$, $\lambda=0.1$\\
relaxation time $1/(2\eta\lambda)$ & steps to forget the initial norm, about $17{,}000$ in that run\\
the timescale $1/(\eta\lambda)$ & the number of recent updates the weights average over
\end{dictionary}

\paragraph{What was thrown away.} For the first idealization, everything except that each step has a fixed length in a norm; convergence will also need a convex loss. For the second, everything about the layer except that the loss ignores its scale; the main theorem also freezes the size of the gradient on the sphere, which in a network changes during training, and assumes steps orthogonal to the weights, which is exact for SGD and false for Adam's preconditioned and momentum steps. Each simplification isolates one mechanism: the ball holds for every sequence of bounded steps, whatever produced them, and the equilibrium needs only orthogonality and a balance of two forces.

\paragraph{What one had to notice.} In \eqref{eq:decay:update} the factor $1-\eta\lambda$ is a convex-combination weight: $(1-\eta\lambda)\theta+\eta\lambda\cdot(-u/\lambda)$ interpolates between the current point and the point $-u/\lambda$, and when the step has a fixed length this is a Frank--Wolfe step. For a scale-invariant layer the gradient step is tangent to the sphere and can only enlarge the norm, while the decay shrinks it; their balance fixes the norm, and the norm fixes how far each step turns the weights.

\begin{fieldguide}{AdamW}{\textcite{LoshchilovHutter2019}}
\fgrow{update}{$m_t=\beta_1m_{t-1}+(1-\beta_1)g_t$, \ $v_t=\beta_2v_{t-1}+(1-\beta_2)g_t^2$, \ $\theta_{t+1}=(1-\eta_t\lambda)\theta_t-\eta_t\,\hat m_t/(\sqrt{\hat v_t}+\epsilon)$}
\fgrow{norm}{$\ell_\infty$ when $\beta_1=\beta_2=0$ (\cref{ch:sign}); with the decay, Frank--Wolfe on $\norm{\theta}_\infty\le1/\lambda$ (\cref{thm:decay:frankwolfe})}
\fgrow{preconditioner}{$\diag(\sqrt{\hat v_t})+\epsilon$ (\cref{ch:adaptive})}
\fgrow{averages}{EMAs of $g$ and $g^2$; through the decay, $\theta_t$ is an EMA of $-u_s/\lambda$ over $1/(\eta\lambda)$ steps (\cref{prop:decay:ema})}
\fgrow{state}{two numbers per parameter}
\fgrow{in the wild}{$\lambda=0.1$, peak $\eta$ from $6\times10^{-5}$ to $6\times10^{-4}$ for language models; matrices decayed, gains and biases usually not}
\fgrow{caution}{codebases differ in what multiplies $\lambda$ (see the notes); Adam with an $L_2$ penalty is a different method (\cref{ex:decay:l2})}
\end{fieldguide}

\begin{fieldguide}{MuonW}{\textcite{JordanEtAl2024Muon}; \textcite{LiuEtAl2025Muon}}
\fgrow{update}{$M_t=\beta M_{t-1}+G_t$, \ $O_t\approx\polar(M_t)$ by Newton--Schulz, \ $W_{t+1}=(1-\eta_t\lambda)W_t-\eta_t\,s\,O_t$}
\fgrow{norm}{spectral (\cref{ch:spectral}); with the decay, Frank--Wolfe on $\norm{W}_{\op}\le s/\lambda$}
\fgrow{averages}{momentum on $G$; the decay averages updates over $1/(\eta\lambda)$ steps}
\fgrow{in the wild}{hidden matrices only, AdamW elsewhere; scale $s$ set to match AdamW's update size; decoupled decay as in AdamW \parencite{LiuEtAl2025Muon}}
\end{fieldguide}

\begin{fieldguide}{Hyperball}{\textcite{WenEtAl2026Hyperball}}
\fgrow{update}{$W_{t+1}=R\,(\hat W_t-\alpha_t\hat U_t)/\norm{\hat W_t-\alpha_t\hat U_t}_F$, with $\hat W=W/R$, $\norm{W_0}_F=R$, and $\hat U_t$ the update of a base optimizer (Adam, Muon) scaled to unit norm}
\fgrow{norm}{the base optimizer's, on the sphere $\norm{W}_F=R$; no decay}
\fgrow{in the wild}{unrefereed 2026 preprint, with reported speedups over weight decay up to $1.2$B parameters (\cref{sec:decay:misses})}
\end{fieldguide}

\section{The theory: decoupled weight decay is Frank--Wolfe}\label{sec:decay:frankwolfe}

To minimize a smooth convex $f$ over a compact convex set $\cK$, the \term{Frank--Wolfe method}\index{Frank--Wolfe method} \parencite{FrankWolfe1956}, or conditional gradient method, minimizes the linearization of $f$ over $\cK$ and moves a fraction $\gamma_t$ of the way to the minimizer:
\[
  v_t\in\argmin_{v\in\cK}\ip{\nabla f(x_t),v},\qquad x_{t+1}=(1-\gamma_t)\,x_t+\gamma_t\,v_t .
\]
It needs only the linear minimization oracle of $\cK$, and its iterates stay in $\cK$ as convex combinations of points of $\cK$. For the ball of a norm the oracle is $\lmo$ rescaled (\cref{fig:decay:pictures}, left). Write $\decayball=\{\theta:\norm{\theta}\le1/\lambda\}$, $\gamma_t=\eta_t\lambda$ and $\Pi_t=\prod_{s<t}(1-\gamma_s)$, with $\Pi_0=1$.

\begin{figure}[t]
\centering
\begin{tikzpicture}[scale=0.92,>=stealth]
  \begin{scope}
    \fill[seamist!70] (-1.6,-1.6) rectangle (1.6,1.6);
    \draw[seadeep,line width=0.7pt] (-1.6,-1.6) rectangle (1.6,1.6);
    \node[font=\decayfont,text=seadeep,anchor=south west] at (-1.6,1.62) {$\norm{\theta}_\infty\le1/\lambda$};
    \coordinate (th) at (0.56,0.08);
    \coordinate (v) at (-1.6,-1.6);
    \coordinate (thn) at ($(th)!0.25!(v)$);
    \draw[muted,densely dashed,line width=0.5pt] (th) -- (v);
    \draw[->,sandline!55!black,line width=0.8pt] (th) -- ++(0.55,0.75) node[anchor=west,inner sep=2pt,font=\decayfont] {$g_t$};
    \draw[->,sea,line width=1.0pt] (th) -- (thn);
    \fill[ink] (th) circle (1.4pt) node[anchor=west,font=\decayfont,xshift=2pt,yshift=-2pt] {$\theta_t$};
    \fill[sea] (thn) circle (1.4pt) node[anchor=north east,font=\decayfont,inner sep=1pt] {$\theta_{t+1}$};
    \fill[seadeep] (v) circle (1.6pt) node[anchor=north,font=\decayfont,yshift=-1pt] {$v_t=-\sign(g_t)/\lambda$};
    \coordinate (a0) at (2.55,1.95);
    \draw[muted,line width=0.5pt,densely dotted] (a0) -- (2.15,1.62) -- (1.82,1.12) -- (1.47,0.95);
    \fill[muted] (a0) circle (1.2pt) node[anchor=west,font=\decayfont] {$\theta_0$};
  \end{scope}
  \begin{scope}[xshift=5.4cm,yshift=-1.4cm]
    \def\R{2.6}
    \draw[muted!70,line width=0.4pt] (0:\R) arc[start angle=0,end angle=62,radius=\R];
    \coordinate (O) at (0,0);
    \coordinate (T) at (\R,0);
    \coordinate (S) at ({0.75*\R},0);
    \coordinate (N) at ({0.75*\R},{sqrt(1-0.75*0.75)*\R});
    \draw[->,ink,line width=0.9pt] (O) -- (T) node[anchor=north,font=\decayfont] {$\theta_t$};
    \draw[->,sandline!55!black,line width=0.9pt] (O) -- (S);
    \node[anchor=north,font=\decayfont,text=sandline!55!black] at ($(S)+(-0.35,-0.04)$) {$(1-\eta\lambda)\theta_t$};
    \draw[->,sea,line width=0.9pt] (S) -- (N) node[midway,anchor=west,font=\decayfont] {$-\eta g_t$};
    \draw[->,seadeep,line width=0.9pt] (O) -- (N) node[anchor=south,font=\decayfont] {$\theta_{t+1}$};
    \draw[muted,line width=0.4pt] ($(S)+(0,0.18)$) -- ++(-0.18,0) -- ++(0,-0.18);
    \draw[muted,line width=0.5pt] (0.75,0) arc[start angle=0,end angle=41.4,radius=0.75];
    \node[font=\decayfont] at (0.98,0.32) {$\omega$};
  \end{scope}
\end{tikzpicture}
\caption{Two pictures of decoupled weight decay, drawn with $\eta\lambda=\frac14$. Left: a step of AdamW without memory moves $\theta_t$ a fraction $\eta\lambda$ of the way to the vertex $-\sign(g_t)/\lambda$ of the ball $\norm{\theta}_\infty\le1/\lambda$, a Frank--Wolfe step; an initialization outside the ball is drawn in (\cref{thm:decay:frankwolfe}). Right: for a scale-invariant layer the gradient is orthogonal to the weights; in equilibrium $\norm{\theta_{t+1}}_2=\norm{\theta_t}_2$, and the right triangle gives $\cos\omega=1-\eta\lambda$ (\cref{thm:decay:equilibrium}).}
\label{fig:decay:pictures}
\end{figure}

\begin{theorem}[Decoupled weight decay is Frank--Wolfe]\label{thm:decay:frankwolfe}
Let the iteration satisfy \cref{ass:decay:nsd}.
\begin{enumerate}
\item With $v_t=\lmo(g_t)/\lambda$ we have $v_t\in\argmin_{v\in\decayball}\ip{g_t,v}$ and $\theta_{t+1}=(1-\gamma_t)\theta_t+\gamma_tv_t$: every step is a Frank--Wolfe step on $\decayball$ with step size $\gamma_t$ for the linear function $\ip{g_t,\cdot}$.
\item For every $t$, $\theta_t-\Pi_t\theta_0\in(1-\Pi_t)\,\decayball$. Consequently the distance from $\theta_t$ to $\decayball$ satisfies
\[
  \bigl(\norm{\theta_t}-1/\lambda\bigr)_+\le\Pi_t\,\bigl(\norm{\theta_0}-1/\lambda\bigr)_+ ,
\]
with equality for every $t$ when $\lmo(g_s)=\theta_0/\norm{\theta_0}$ for all $s$. So $\decayball$ is invariant and, at a constant step size, attracts every initialization at the rate $(1-\eta\lambda)^t$.
\item Let $\Loss$ be convex and $L$-smooth with respect to $\norm{\cdot}$, let $g_t=\nabla\Loss(\theta_t)$, and let $h_t=\Loss(\theta_t)-\Loss^\star_\lambda$ with $\Loss^\star_\lambda=\min_{\decayball}\Loss$. Then for every $t$
\begin{equation}\label{eq:decay:fwrec}
  h_{t+1}\le(1-\gamma_t)\,h_t+\tfrac L2\,\gamma_t^2\,\norm{v_t-\theta_t}^2 .
\end{equation}
For $\gamma_t=2/(t+2)$, that is $\eta_t=2/(\lambda(t+2))$, and $D=1/\lambda+\max\{1/\lambda,\norm{\theta_0}\}$, we have $h_t\le2LD^2/(t+2)$ for $t\ge1$, which is $8L/(\lambda^2(t+2))$ if $\theta_0\in\decayball$. For a constant $\gamma_t=\gamma$, $\limsup_th_t\le2L\gamma/\lambda^2=2L\eta/\lambda$, and $h_t\le(1-\gamma)^th_0+2L\gamma/\lambda^2$ for every $t$ if $\theta_0\in\decayball$.
\item If instead $g_t=\nabla\Loss(\theta_t)+\zeta_t$ with an error $\zeta_t$, then \eqref{eq:decay:fwrec} holds with the additional term $2\gamma_t\norm{\zeta_t}_*/\lambda$ on its right side.
\end{enumerate}
\end{theorem}

\begin{proof}
The plan is to read (i) off the definition of $\lmo$, to unroll the recursion for (ii), and to repeat the classical Frank--Wolfe argument for (iii) and (iv), keeping track of iterates that start outside the ball.

(i) The map $u\mapsto u/\lambda$ is a bijection from the unit ball onto $\decayball$ and multiplies $\ip{g_t,u}$ by $1/\lambda>0$, so it maps the minimizers of $\ip{g_t,\cdot}$ over the unit ball onto those over $\decayball$. The update is the identity $\eta_t\lmo(g_t)=\eta_t\lambda\cdot\lmo(g_t)/\lambda=\gamma_tv_t$.

(ii) Write $\Pi_{s,t}=\prod_{s\le k<t}(1-\gamma_k)$, so that $\Pi_{0,t}=\Pi_t$ and $\Pi_{t,t}=1$. By induction on $t$, (i) gives
\[
  \theta_t=\Pi_t\theta_0+\sum_{s<t}w_s\,v_s,\qquad w_s=\gamma_s\Pi_{s+1,t} .
\]
The weights are nonnegative because $0<\gamma_k\le1$, and they telescope, $w_s=\Pi_{s+1,t}-\Pi_{s,t}$, to $1-\Pi_t$. If $\Pi_t<1$, then $\sum_sw_sv_s/(1-\Pi_t)$ is a convex combination of points of $\decayball$, so $\sum_sw_sv_s\in(1-\Pi_t)\decayball$; if $\Pi_t=1$ the sum is zero. The triangle inequality gives $\norm{\theta_t}-1/\lambda\le\Pi_t(\norm{\theta_0}-1/\lambda)$, and we take positive parts, as $\Pi_t\ge0$. If $\lmo(g_s)=\theta_0/\norm{\theta_0}$ for all $s$, every $v_s$ is a positive multiple of $\theta_0$ and the triangle inequality is an equality. Finally, in every norm the distance from $\theta$ to $\decayball$ is $(\norm{\theta}-1/\lambda)_+$: if $\norm{\theta}>1/\lambda$, the point $\theta/(\lambda\norm{\theta})$ of the ball is at distance $\norm{\theta}-1/\lambda$, and every $x\in\decayball$ has $\norm{\theta-x}\ge\norm{\theta}-\norm{x}\ge\norm{\theta}-1/\lambda$.

(iii) The minimum $\Loss^\star_\lambda$ is attained at some $\theta^\star\in\decayball$, since $\decayball$ is compact. The descent lemma (\cref{lem:toolkit:descent}), applied to the step $\gamma_t(v_t-\theta_t)$, gives
\[
  \Loss(\theta_{t+1})\le\Loss(\theta_t)+\gamma_t\ip{\nabla\Loss(\theta_t),v_t-\theta_t}+\tfrac L2\gamma_t^2\norm{v_t-\theta_t}^2 .
\]
Because $v_t$ minimizes $\ip{\nabla\Loss(\theta_t),\cdot}$ over $\decayball\ni\theta^\star$, and $\Loss$ is convex,
\[
  \ip{\nabla\Loss(\theta_t),v_t-\theta_t}\le\ip{\nabla\Loss(\theta_t),\theta^\star-\theta_t}\le\Loss(\theta^\star)-\Loss(\theta_t)=-h_t .
\]
Substituting gives \eqref{eq:decay:fwrec}, which does not use the sign of $h_t$ (negative values are possible while $\theta_t$ is outside the ball). By (ii), $\norm{v_t-\theta_t}\le1/\lambda+\norm{\theta_t}\le2/\lambda+\Pi_te_0$ with $e_0=(\norm{\theta_0}-1/\lambda)_+$.

For $\gamma_t=2/(t+2)$ we have $\gamma_0=1$, so the first step lands on the vertex $v_0$ and $\Pi_t=0$ for $t\ge1$; hence $\norm{v_t-\theta_t}\le2/\lambda\le D$ for $t\ge1$, and $\norm{v_0-\theta_0}\le1/\lambda+\norm{\theta_0}\le D$. We show $h_t\le2LD^2/(t+2)$ by induction. At $t=1$, \eqref{eq:decay:fwrec} with $\gamma_0=1$ gives $h_1\le\frac L2D^2\le\frac23LD^2$. If the bound holds at $t\ge1$, then
\[
  h_{t+1}\le\frac t{t+2}\cdot\frac{2LD^2}{t+2}+\frac L2\cdot\frac4{(t+2)^2}\,D^2=\frac{2LD^2(t+1)}{(t+2)^2}\le\frac{2LD^2}{t+3},
\]
because $(t+1)(t+3)\le(t+2)^2$. If $\theta_0\in\decayball$, then $D=2/\lambda$.

For a constant $\gamma$, unrolling \eqref{eq:decay:fwrec} gives
\[
  h_t\le(1-\gamma)^th_0+\frac L2\gamma^2\sum_{s<t}(1-\gamma)^{t-1-s}\Bigl(\frac2\lambda+(1-\gamma)^se_0\Bigr)^2 .
\]
Expanding the square, the sum is at most $4/(\lambda^2\gamma)+(4e_0/\lambda)\,t(1-\gamma)^{t-1}+e_0^2(1-\gamma)^{t-1}/\gamma$, whose last two terms tend to zero; if $\theta_0\in\decayball$, they vanish.

(iv) Now $v_t$ minimizes $\ip{g_t,\cdot}$, so $\ip{g_t,v_t}\le\ip{g_t,\theta^\star}$, and with $\nabla\Loss(\theta_t)=g_t-\zeta_t$,
\[
  \ip{\nabla\Loss(\theta_t),v_t-\theta_t}\le\ip{\nabla\Loss(\theta_t),\theta^\star-\theta_t}+\ip{\zeta_t,\theta^\star-v_t} .
\]
The first term is at most $-h_t$ as in (iii); the second is at most $\norm{\zeta_t}_*\norm{\theta^\star-v_t}\le2\norm{\zeta_t}_*/\lambda$ by Hölder's inequality (\cref{lem:toolkit:dualnorm}). The rest of (iii) is unchanged.
\end{proof}

So the schedule $\eta_t$ is a Frank--Wolfe step-size schedule, $\gamma_t=\eta_t\lambda$, on a ball whose radius $1/\lambda$ does not depend on the schedule.\sidenote{The radius depends on what multiplies $\lambda$. In PyTorch the decay per step is $\eta_t\lambda$ and the radius $1/\lambda$; in the AdamW of \textcite{LoshchilovHutter2019} the decay is the schedule multiplier times $\lambda$ while the step also carries the base learning rate $\alpha$, and the radius is $\alpha/\lambda$.} The classical $\gamma_t=2/(t+2)$ is a learning rate that starts at $1/\lambda$ and decays like $1/t$; every schedule with $\gamma_t\to0$ and $\sum_t\gamma_t=\infty$ drives $h_t$ to zero for exact gradients (\cref{ex:decay:schedules}), and a constant one leaves a neighborhood of size at most $2L\eta/\lambda$, which a decaying schedule removes, as in \cref{ch:schedules}. These are worst cases: on the quadratics of \cref{fig:decay:frankwolfe} the gap falls like $1/t^2$ under $2/(t+2)$ and is of order $\gamma^2$ at a constant step (\cref{ex:decay:cycle}). Part (iv) adds a floor that no schedule removes: if $\E\norm{\zeta_t}_*\le\sigma_*$, the bound on $\E h_t$ does not fall below $2\sigma_*/\lambda$. The error can be that of a stochastic gradient or of a momentum average, which \cref{ch:momentum} shows how to make small; this is why Lion and Muon feed a momentum to their oracle.

Three optimizers of this book are instances. AdamW without memory is Frank--Wolfe on the $\ell_\infty$ ball, and for AdamW with memory \textcite{XieLi2024} proved that if the iterates converge under a nonincreasing schedule with divergent sum, the limit is a Karush--Kuhn--Tucker point of $\min\Loss$ subject to $\norm{\theta}_\infty\le1/\lambda$. \Textcite{ChenEtAl2024Lion} proved the corresponding statement for Lion with a Lyapunov function, Muon with decay solves a problem constrained in the spectral norm \parencite{ChenLiLiu2025}, and \textcite{PethickEtAl2025} built optimizers directly as conditional-gradient methods with the oracles of \cref{ch:steepest}. Bounded steps that are not normalized are drawn in the same way.

\begin{proposition}[Bounded steps]\label[proposition]{prop:decay:bounded}
If $\theta_{t+1}=(1-\eta_t\lambda)\theta_t-\eta_tu_t$ with $0<\eta_t\lambda\le1$ and $\norm{u_t}\le C$ for every $t$, then $\theta_t-\Pi_t\theta_0\in(1-\Pi_t)\,C\decayball$ and $(\norm{\theta_t}-C/\lambda)_+\le\Pi_t(\norm{\theta_0}-C/\lambda)_+$.
\end{proposition}

\begin{proof}
Put $v_t=-u_t/\lambda\in C\decayball$ and repeat the proof of \cref{thm:decay:frankwolfe}(ii) with $C\decayball$ in place of $\decayball$.
\end{proof}

Adam's step is bounded coordinatewise by $(1-\beta_1)/\sqrt{(1-\beta_2)(1-\beta_1^2/\beta_2)}$ when $\beta_1^2<\beta_2$ (\cref{ch:sign}), which is $1.17$ for $(\beta_1,\beta_2)=(0.9,0.95)$ and $7.27$ for $(0.9,0.999)$, so AdamW's weights are drawn into the $\ell_\infty$ ball of radius $1.17/\lambda$ or $7.27/\lambda$.\sidenote{The bound is proved in \cref{ch:sign} without bias correction. With bias correction the Cauchy--Schwarz bound after $t$ steps increases with $t$ to the same limit, as we checked numerically for these pairs.} Where in the ball do the weights sit? One coordinate of sign descent answers exactly.

\begin{proposition}[Consistent and noisy coordinates]\label[proposition]{prop:decay:regimes}
Let $0<\gamma=\eta\lambda<1$ and $x_{t+1}=(1-\gamma)x_t-\eta s_t$, where $\abs{s_t}\le1$ and $\E[s_t\mid\cF_t]=\mu$ for every $t$, with $\cF_t$ generated by $x_0,s_0,\dots,s_{t-1}$.
\begin{enumerate}
\item $\E x_t=-\mu/\lambda+(1-\gamma)^t(x_0+\mu/\lambda)$; if $s_t=\mu\in\{-1,1\}$ is deterministic, $x_t$ follows this curve to the vertex $-\mu/\lambda$.
\item If $\mu=0$ and $\E[s_t^2\mid\cF_t]=1$, as for random signs, then $\E x_t^2=(1-\gamma)^{2t}x_0^2+\eta^2\bigl(1-(1-\gamma)^{2t}\bigr)/\bigl(1-(1-\gamma)^2\bigr)$, which tends to $\eta/(\lambda(2-\eta\lambda))$: the coordinate settles at root mean square $\lambda^{-1}\sqrt{\gamma/(2-\gamma)}\approx\lambda^{-1}\sqrt{\eta\lambda/2}$.
\end{enumerate}
\end{proposition}

\begin{proof}
Taking expectations, $\E x_{t+1}=(1-\gamma)\E x_t-\eta\mu$, whose solution is (i). For (ii), $\E[x_ts_t]=\E[x_t\,\E[s_t\mid\cF_t]]=0$ because $x_t$ is determined by $\cF_t$, so $\E x_{t+1}^2=(1-\gamma)^2\E x_t^2+\eta^2$.
\end{proof}

For sign descent on a coordinate whose gradient has mean $\bar g_i$ and Gaussian noise of standard deviation $\sigma_i$, the step has mean $\mu=\erf(\bar g_i/(\sigma_i\sqrt2))$, so the coordinate settles at $\erf(\bar g_i/(\sigma_i\sqrt2))/\lambda$ (with the sign that opposes the gradient): at the vertex when the gradient is consistent, deep inside the ball when it is noise. For the GPT-3 models $\eta\lambda$ is between $6\times10^{-6}$ and $6\times10^{-5}$, and a coordinate whose gradient is noise sits at $0.2\%$ to $0.5\%$ of the radius. The ball is a constraint only for coordinates whose updates keep one sign for about $1/(\eta\lambda)$ steps; for the others, the decay sets a size by balancing noise against shrinkage. For scale-invariant weights this balance is the whole story.

\section{The theory: scale invariance and the rotational equilibrium}\label{sec:decay:equilibrium}

Gradients in this section are Euclidean. Two facts make scale-invariant weights special.

\begin{lemma}[Scale invariance]\label[lemma]{lem:decay:scaleinv}
Let $\Loss$ be scale invariant and twice differentiable on $\R^d\setminus\{0\}$. For every $\theta\ne0$ and $c>0$:
(i) $\ip{\nabla\Loss(\theta),\theta}=0$;
(ii) $\nabla\Loss(c\theta)=c^{-1}\nabla\Loss(\theta)$, so $\nabla\Loss(\theta)=\nabla\Loss(\hat\theta)/\norm{\theta}_2$;
(iii) $\nabla^2\Loss(c\theta)=c^{-2}\nabla^2\Loss(\theta)$, so the sharpness is $S(\theta)=S(\hat\theta)/\norm{\theta}_2^2$.
\end{lemma}

\begin{proof}
Differentiate $\Loss(c\theta)=\Loss(\theta)$ in $c$ at $c=1$ for (i), and in $\theta$, once and twice, for $c\nabla\Loss(c\theta)=\nabla\Loss(\theta)$ and $c^2\nabla^2\Loss(c\theta)=\nabla^2\Loss(\theta)$; then take $c=\norm{\theta}_2$ and replace $\theta$ by $\hat\theta$.
\end{proof}

A gradient step is therefore tangent to the sphere, and gradient and curvature fall as the norm grows. By (iii), the stability threshold $\eta S<2$ of \cref{ch:eos} reads $(\eta/\norm{\theta}_2^2)\,S(\hat\theta)<2$: the step size that matters is $\eta/\norm{\theta}_2^2$. If each per-example loss is scale invariant, (i) holds for every stochastic gradient, and one step of \eqref{eq:decay:update} is a right triangle (\cref{fig:decay:pictures}, right).

\begin{lemma}[One step]\label[lemma]{lem:decay:step}
Let $\theta\ne0$, $0<\eta\lambda<1$, $\theta'=(1-\eta\lambda)\theta-\eta u$ with $\ip{u,\theta}=0$, and let $\omega$ be the angle between $\theta$ and $\theta'$. Then
\begin{enumerate}
\item $\norm{\theta'}_2^2=(1-\eta\lambda)^2\norm{\theta}_2^2+\eta^2\norm{u}_2^2$;
\item $\cos\omega=(1-\eta\lambda)\norm{\theta}_2/\norm{\theta'}_2$ and $\tan\omega=\eta\norm{u}_2/((1-\eta\lambda)\norm{\theta}_2)$;
\item if $u\ne0$, then $\hat\theta'=(\hat\theta-\alpha\hat u)/\norm{\hat\theta-\alpha\hat u}_2$ with $\hat u=u/\norm{u}_2$ and $\alpha=\tan\omega$;
\item $\norm{\theta'-\theta}_2^2/\norm{\theta}_2^2=\eta^2\lambda^2+\eta^2\norm{u}_2^2/\norm{\theta}_2^2$.
\end{enumerate}
\end{lemma}

\begin{proof}
Pythagoras' theorem for the orthogonal vectors $\theta$ and $u$, with $\ip{\theta,\theta'}=(1-\eta\lambda)\norm{\theta}_2^2$ for (ii); for (iii), divide $\theta'$ by $(1-\eta\lambda)\norm{\theta}_2>0$; for (iv), $\theta'-\theta=-\eta\lambda\theta-\eta u$.
\end{proof}

Write $a=(1-\eta\lambda)^2$; then $\cos\omega=\sqrt a$ exactly when the norm does not change. The weakest form of the equilibrium needs nothing more.

\begin{proposition}[The equilibrium angle is a time average]\label[proposition]{prop:decay:average}
Let $\theta_{t+1}=(1-\eta\lambda)\theta_t-\eta u_t$ with $\ip{u_t,\theta_t}=0$, $\theta_t\ne0$, constant $\eta$ and $\lambda$, and $q_t=\norm{\theta_t}_2^2$. For all $T_1<T_2$,
\[
  \sum_{t=T_1}^{T_2-1}q_t\tan^2\omega_t=\frac{1-a}a\sum_{t=T_1}^{T_2-1}q_t+\frac{q_{T_2}-q_{T_1}}a .
\]
If $q_{T_2}=q_{T_1}$, the $q$-weighted mean of $\tan^2\omega_t$ over the window is exactly $(1-a)/a=\tan^2\bigl(\arccos(1-\eta\lambda)\bigr)$, whatever the sizes of the steps.
\end{proposition}

\begin{proof}
By \cref{lem:decay:step}(i) and (ii), $aq_t\tan^2\omega_t=\eta^2\norm{u_t}_2^2=q_{t+1}-aq_t=(q_{t+1}-q_t)+(1-a)q_t$. Sum over the window and divide by $a$.
\end{proof}

This is why measured angles are robust: gradients may be large in some steps and small in others, but once decay and growth cancel on average, the mean rotation is fixed by $\eta\lambda$ alone. With a constant gradient size on the sphere, the norm is driven to that balance from every start.

\begin{theorem}[Rotational equilibrium of SGD]\label{thm:decay:equilibrium}
Let $\theta_0\ne0$, $0<\eta\lambda<1$, and $\theta_{t+1}=(1-\eta\lambda)\theta_t-\eta g_t$ with $\ip{g_t,\theta_t}=0$ and $\norm{g_t}_2\norm{\theta_t}_2=\Gamma>0$ for every $t$. Let $q_t=\norm{\theta_t}_2^2$ and
\begin{equation}\label{eq:decay:qinf}
  q_\infty=\frac{\eta\Gamma}{\sqrt{1-a}}=\Gamma\sqrt{\frac{\eta}{\lambda(2-\eta\lambda)}}\approx\Gamma\sqrt{\frac{\eta}{2\lambda}} .
\end{equation}
\begin{enumerate}
\item $q_{t+1}=aq_t+\eta^2\Gamma^2/q_t$, and $q_t\to q_\infty$ for every $\theta_0\ne0$; the distance $\abs{\log(q_t/q_\infty)}$ never increases, and decreases strictly unless $q_t=q_\infty$.
\item Suppose $\eta\lambda\le1-1/\sqrt2$, so that $a\ge\frac12$. If $q_t\ge q_\infty$, then $0\le q_{t+1}-q_\infty\le a\,(q_t-q_\infty)$. If $q_t<q_\infty$, then $q_{t+1}>q_t$, and either $q_{t+1}\ge q_\infty$ or $0<q_\infty-q_{t+1}\le(2a-1)(q_\infty-q_t)$. Near $q_\infty$ the error is multiplied by $2a-1=1-4\eta\lambda+2\eta^2\lambda^2$ per step.
\item As $t\to\infty$: the angle $\omega_t$ between $\theta_t$ and $\theta_{t+1}$ tends to $\omega_\infty=\arccos(1-\eta\lambda)=\sqrt{2\eta\lambda}\,(1+O(\eta\lambda))$; the relative step $\norm{\theta_{t+1}-\theta_t}_2/\norm{\theta_t}_2$ tends to $\sqrt{2\eta\lambda}$ exactly; and the direction takes gradient steps on the unit sphere, $\hat\theta_{t+1}=(\hat\theta_t-\eta_{\mathrm{eff},t}\hat g_t)/\norm{\hat\theta_t-\eta_{\mathrm{eff},t}\hat g_t}_2$ with $\hat g_t=\norm{\theta_t}_2\,g_t$ the gradient at $\hat\theta_t$, whose \term{effective step size}\index{effective learning rate} $\eta_{\mathrm{eff},t}=\eta/((1-\eta\lambda)q_t)$ tends to $\tan(\omega_\infty)/\Gamma$.
\item If instead $\Gamma_t=\norm{g_t}_2\norm{\theta_t}_2\in[\Gamma_1,\Gamma_2]$ for every $t$, with $a\ge\frac12$ and $\Gamma_1/\Gamma_2>(1-a)/a$, then $q_\infty(\Gamma_1)\le\liminf_tq_t\le\limsup_tq_t\le q_\infty(\Gamma_2)$, where $q_\infty(\Gamma)$ is given by \eqref{eq:decay:qinf}.
\end{enumerate}
\end{theorem}

\begin{proof}
The plan is to reduce the dynamics to one scalar map, to show that the map moves every point toward its fixed point on a logarithmic scale, and to read the angles off \cref{lem:decay:step}.

(i) \Cref{lem:decay:step}(i) with $\norm{g_t}_2^2=\Gamma^2/q_t$ gives the recursion, and $q_{t+1}\ge aq_t>0$. Put $x_t=q_t/q_\infty$. Since $\eta^2\Gamma^2=(1-a)q_\infty^2$, dividing by $q_\infty$ gives $x_{t+1}=\varphi(x_t)$ with
\[
  \varphi(x)=a\,x+(1-a)\,\frac1x ,
\]
a convex combination, with both weights in $(0,1)$, of $x$ and $1/x$. If $x\ne1$, these two numbers lie on opposite sides of $1$ and $\varphi(x)$ lies strictly between them, so $\abs{\log\varphi(x)}<\abs{\log x}$. The numbers $\ell_t=\abs{\log x_t}$ therefore decrease to a limit $\ell\ge0$. Suppose $\ell>0$. The continuous function $y\mapsto\abs y-\abs{\log\varphi(e^y)}$ is positive on the compact set $\{y:\ell\le\abs y\le\ell_0\}$, which contains every $\log x_t$, so it has a positive minimum $\delta$ there, and $\ell_t\le\ell_0-t\delta$, which contradicts $\ell_t\ge\ell$ for $t>\ell_0/\delta$. Hence $\ell=0$.

(ii) Factor the distances:
\begin{equation}\label{eq:decay:factor}
  \varphi(x)-1=(x-1)\Bigl(a-\frac{1-a}x\Bigr),\qquad \varphi(x)-x=(1-a)\,\frac{1-x^2}x .
\end{equation}
If $x\ge1$, the second factor of the first identity lies in $[2a-1,a)\subseteq[0,a)$ because $a\ge\frac12$, so $0\le\varphi(x)-1\le a(x-1)$. If $x<1$, the second identity gives $\varphi(x)>x$, and the second factor of the first is less than $2a-1$: either it is negative and $\varphi(x)>1$, or $0\le1-\varphi(x)<(2a-1)(1-x)$, where $\varphi(x)=1$ is impossible because $\varphi(x)$ lies strictly between $x$ and $1/x$. Multiply by $q_\infty$. The local factor is $\varphi'(1)=2a-1$.

(iii) Since $q_{t+1}/q_t\to1$, \cref{lem:decay:step}(ii) gives $\cos\omega_t=(1-\eta\lambda)\sqrt{q_t/q_{t+1}}\to1-\eta\lambda$, and \cref{lem:decay:step}(iv) gives $\norm{\theta_{t+1}-\theta_t}_2^2/q_t=\eta^2\lambda^2+\eta^2\Gamma^2/q_t^2\to\eta^2\lambda^2+(1-a)=2\eta\lambda$; the expansion $\arccos(1-z)=\sqrt{2z}\,(1+z/12+O(z^2))$ gives the size of $\omega_\infty$. For the direction, \cref{lem:decay:scaleinv}(ii) for the per-example loss gives $g_t=\hat g_t/\norm{\theta_t}_2$, and \cref{lem:decay:step}(iii) gives the update with $\alpha_t\hat u_t=\eta g_t/((1-\eta\lambda)\norm{\theta_t}_2)=\eta_{\mathrm{eff},t}\hat g_t$. The limit of $\eta_{\mathrm{eff},t}$ is $\eta/((1-\eta\lambda)q_\infty)=\sqrt{1-a}/(\sqrt a\,\Gamma)$, and $\sqrt{1-a}/\sqrt a=\tan\omega_\infty$ because $\cos\omega_\infty=\sqrt a$.

(iv) Let $b_t=\eta^2\Gamma_t^2$, $b_i=\eta^2\Gamma_i^2$ and $q_i=q_\infty(\Gamma_i)=\sqrt{b_i/(1-a)}$, so that $q_{t+1}=aq_t+b_t/q_t$. \emph{Lower bound.} If $q_t<q_1$, then $q_{t+1}-q_t\ge b_1/q_t-(1-a)q_t=(1-a)(q_1^2-q_t^2)/q_t\ge(1-a)(q_1-q_t)$, so $q_1-q_{t+1}\le a(q_1-q_t)$. If $q_t\ge q_1$, then $q_{t+1}\ge aq_t+b_1/q_t\ge q_1$ by (ii) with $\Gamma=\Gamma_1$. Hence $(q_1-q_t)_+\le a^t(q_1-q_0)_+$. \emph{Upper bound.} If $q_t\ge q_2$, then $q_{t+1}\le aq_t+b_2/q_2=aq_t+(1-a)q_2$, so $q_{t+1}-q_2\le a(q_t-q_2)$. If $q_t<q_2$, then $q_{t+1}\le aq_t+b_2/q_t$, which by \eqref{eq:decay:factor} with $x=q_t/q_2$ is at most $q_2$ once $q_t\ge(1-a)q_2/a$. Since $q_1/q_2=\Gamma_1/\Gamma_2>(1-a)/a$, the lower bound gives $q_t\ge(1-a)q_2/a$ for all large $t$, and from then on $(q_{t+1}-q_2)_+\le a(q_t-q_2)_+$.
\end{proof}

The theorem corrects a natural guess about time scales. Far above equilibrium the squared norm decays like $a^t\approx e^{-2\eta\lambda t}$, as if the gradient were absent; near equilibrium it relaxes at the rate $4\eta\lambda$, twice as fast, because a smaller norm means a larger gradient (\cref{lem:decay:scaleinv}(ii)). From far below, it rises like $\sqrt{2(1-a)t}\,q_\infty$ and arrives within about $1/(4\eta\lambda)$ steps, whatever its starting value. Optimizers whose step has a fixed size, such as AdamW, Lion and Muon, have no such restoring force, and their recursion is linear.

\begin{proposition}[Steps of fixed size: AdamW, Lion, Muon]\label[proposition]{prop:decay:adamw}
Let $\theta_{t+1}=(1-\eta_t\lambda)\theta_t-\eta_tu_t$ with $0<\eta_t\lambda<1$ and $q_t=\norm{\theta_t}_2^2$.
\begin{enumerate}
\item If $\ip{u_t,\theta_t}=0$ and $\norm{u_t}_2^2=\rho^2d$ for every $t$, then $q_{t+1}=(1-\eta_t\lambda)^2q_t+\eta_t^2\rho^2d$ exactly. At a constant step size, $q_t-q_\infty=a^t(q_0-q_\infty)$ with $q_\infty=\eta^2\rho^2d/(1-a)=\rho^2\eta d/(\lambda(2-\eta\lambda))$, and $\cos\omega_t\to1-\eta\lambda$.
\item If instead $\E[\ip{u_t,\theta_t}\mid\cF_t]=0$ and $\E[\norm{u_t}_2^2\mid\cF_t]=\rho^2d$, with $\cF_t$ the past up to $\theta_t$, the same recursion holds for $\E q_t$.
\end{enumerate}
\end{proposition}

\begin{proof}
Part (i) is \cref{lem:decay:step}(i) and (ii) and the solution of a linear recursion. For (ii), expand $q_{t+1}=(1-\eta_t\lambda)^2q_t-2\eta_t(1-\eta_t\lambda)\ip{u_t,\theta_t}+\eta_t^2\norm{u_t}_2^2$ and take conditional expectations.
\end{proof}

With steps of root-mean-square size $\rho$ per entry, the equilibrium entries have root mean square $\rho\sqrt{\eta/(\lambda(2-\eta\lambda))}\approx\rho\sqrt{\eta/(2\lambda)}$, the noisy-coordinate value of \cref{prop:decay:regimes}, and the angle is that of SGD. Sign steps have $\rho=1$; an orthogonalized $d_{\mathrm{out}}\times d_{\mathrm{in}}$ step scaled by $s$ has $\rho^2d=s^2\min\{d_{\mathrm{out}},d_{\mathrm{in}}\}$. Orthogonality holds exactly for normalized gradient steps $u_t=\rho\sqrt d\,g_t/\norm{g_t}_2$; for sign steps it holds on average when the gradient on the sphere is symmetric noise, which is the hypothesis of (ii); for preconditioned steps and for momentum it fails (\cref{sec:decay:reading}). For momentum we have a derivation rather than a proof.

\begin{heuristic}
\emph{Momentum and the angle per step.} Model AdamW's step, when the gradient is noise, by $u_t=m_t=\beta_1m_{t-1}+(1-\beta_1)\xi_t$ with independent $\xi_t$ of mean zero and $\E\norm{\xi_t}_2^2=d$, the role of $g/\sqrt{\hat v}$. Then $\E\norm{u_t}_2^2\to d(1-\beta_1)/(1+\beta_1)$, while a sum of $k\gg1/(1-\beta_1)$ consecutive steps has second moment about $kd$, because the weights $(1-\beta_1)\beta_1^j$ sum to one: momentum smooths the walk without slowing its diffusion. If $1/(\eta\lambda)\gg1/(1-\beta_1)$, the weights average about $1/(\eta\lambda)$ steps, so their squared norm keeps the value $\eta d/(2\lambda)$ of \cref{prop:decay:adamw} (\cref{ex:decay:momentum} computes the correction), and the angle per step is
\begin{equation}\label{eq:decay:momentum}
  \omega\approx\sqrt{\frac{\eta^2\E\norm{u_t}_2^2}{\eta d/(2\lambda)}}=\sqrt{2\eta\lambda\,\frac{1-\beta_1}{1+\beta_1}},
\end{equation}
a factor $0.23$ below $\sqrt{2\eta\lambda}$ for $\beta_1=0.9$, while over $k$ steps the direction still turns by about $\sqrt{2\eta\lambda k}$. In the balance of \cref{prop:decay:average}, the rest of the decay is carried by the component of $m_t$ along $\theta_t$: the averaged past gradients were tangent to past directions and point partly outward from the present one.
\end{heuristic}

The equilibrium moves when the step size does, and the norm follows with a lag.

\begin{proposition}[The rotation lags the schedule]\label[proposition]{prop:decay:schedule}
In \cref{prop:decay:adamw}(i), let $q_\infty(\eta)=\rho^2\eta d/(\lambda(2-\eta\lambda))$ and $\omega_\infty(\eta)=\arccos(1-\eta\lambda)$. If $\eta_t$ is nonincreasing and $q_0\ge q_\infty(\eta_0)$, then $q_t\ge q_\infty(\eta_t)$ and $\omega_t\le\omega_\infty(\eta_t)$ for every $t$. If $\eta_t$ is nondecreasing and $q_0\le q_\infty(\eta_0)$, both inequalities reverse.
\end{proposition}

\begin{proof}
The function $q_\infty$ increases on $(0,1/\lambda)$, and $q_\infty(\eta_t)$ is the fixed point of the increasing affine map $q\mapsto(1-\eta_t\lambda)^2q+\eta_t^2\rho^2d$. So $q_t\ge q_\infty(\eta_t)$ implies $q_{t+1}\ge q_\infty(\eta_t)\ge q_\infty(\eta_{t+1})$, and induction proves the first claim. By \cref{lem:decay:step}(ii), $\tan\omega_t=\eta_t\rho\sqrt d/((1-\eta_t\lambda)\sqrt{q_t})$, which is at most the same expression with $q_\infty(\eta_t)$ in place of $q_t$, namely $\tan\omega_\infty(\eta_t)$. The nondecreasing case reverses every inequality.
\end{proof}

A decaying schedule therefore cools a scale-invariant layer faster than the quasi-static value $\sqrt{2\eta_t\lambda}$ suggests: the relaxation time $1/(2\eta_t\lambda)$ grows as $\eta_t$ falls, the norm cannot shrink in time, and near the end of a decay to zero the rotation falls like $\eta_t$ rather than $\sqrt{\eta_t}$ (\cref{ex:decay:lag}). A warmup that starts from a small norm rotates the weights faster than its equilibrium value. Because only the direction matters, all of this can be said with a learning rate on the sphere.

\begin{proposition}[Weight decay is a Hyperball schedule]\label[proposition]{prop:decay:hyperball}
Let $\Loss$ and each $\ell(\cdot;\xi)$ be scale invariant, and let the base step be $u_t=U(\hat\theta_t,\xi_t)$, a function of the direction and of the randomness $\xi_t$ only, with $\ip{u_t,\theta_t}=0$ and $\norm{u_t}_2^2=\rho^2d$. Let $q_t$ solve the recursion of \cref{prop:decay:adamw}(i) from $q_0=\norm{\theta_0}_2^2$, and run Hyperball on the unit sphere, $\hat w_{t+1}=(\hat w_t-\alpha_t\hat U_t)/\norm{\hat w_t-\alpha_t\hat U_t}_2$ with $\hat U_t=U(\hat w_t,\xi_t)/\norm{U(\hat w_t,\xi_t)}_2$, from $\hat w_0=\hat\theta_0$, with the same $\xi_t$ and
\[
  \alpha_t=\frac{\eta_t\rho\sqrt d}{(1-\eta_t\lambda)\sqrt{q_t}} .
\]
Then $\hat w_t=\hat\theta_t$ for every $t$, and the two runs have the same losses. At a constant step size, $\alpha_t\to\tan\omega_\infty\approx\sqrt{2\eta\lambda}$.
\end{proposition}

\begin{proof}
By \cref{prop:decay:adamw}(i), $\norm{\theta_t}_2^2=q_t$. If $\hat w_t=\hat\theta_t$, then $\hat U_t=\hat u_t$, and \cref{lem:decay:step}(iii) gives $\hat\theta_{t+1}=\hat w_{t+1}$. Induct; scale invariance gives equal losses; the limit follows from $q_t\to q_\infty$.
\end{proof}

The hypotheses hold exactly for normalized gradient steps, whose direction depends on $\theta$ only through $\hat\theta$ by \cref{lem:decay:scaleinv}(ii); for sign, orthogonalized and momentum steps they hold approximately, up to a radial component that averages out and a norm that changes slowly. The last result looks at the weights themselves.

\begin{proposition}[AdamW is an exponential moving average of its updates]\label[proposition]{prop:decay:ema}
For \eqref{eq:decay:update} with $0<\eta_t\lambda\le1$ and $\Pi_{s,t}=\prod_{s\le k<t}(1-\eta_k\lambda)$,
\[
  \theta_t=\Pi_{0,t}\,\theta_0-\sum_{s<t}\eta_s\,\Pi_{s+1,t}\,u_s ,\qquad \Pi_{s,t}\le\exp\Bigl(-\lambda\sum_{s\le k<t}\eta_k\Bigr).
\]
At a constant step size, $\theta_t=(1-\eta\lambda)^t\theta_0+\sum_{s<t}\eta\lambda(1-\eta\lambda)^{t-1-s}(-u_s/\lambda)$, an exponential moving average of the points $-u_s/\lambda$ with timescale $\tau=1/(\eta\lambda)$ steps, in which the updates older than $k$ steps carry total weight at most $(1-\eta\lambda)^k$.
\end{proposition}

\begin{proof}
Induction on $t$; $1-x\le e^{-x}$; and $\sum_{s<t-k}\eta\lambda(1-\eta\lambda)^{t-1-s}\le(1-\eta\lambda)^k$.
\end{proof}

The proposition is \cref{thm:decay:frankwolfe}(ii) read as an average, and with one hypothesis it gives scaling rules. A run of $T$ steps covers $\lambda\sum_t\eta_t$ timescales, $\eta\lambda T$ at a constant step size. If the right number of timescales is a fixed fraction of the run, as the evidence of \textcite{WangAitchison2024} suggests (\cref{obs:decay:horizon}), then $\eta\lambda T$ should be held fixed: at fixed learning rate and batch size, $\lambda$ should fall like the inverse of the number of training tokens, not like its inverse square root as in the rule of \textcite{LoshchilovHutter2019}. When $\mu$P scales AdamW's hidden-layer learning rate like $1/d_{\mathrm{in}}$ (\cref{ch:modular}), holding $\eta\lambda$ fixed requires $\lambda\propto d_{\mathrm{in}}$; with a fixed $\lambda$, wider models rotate more slowly and reach equilibrium later. Both rules follow from the identity and the hypothesis; neither is proved optimal.

\section{Reading the theory}\label{sec:decay:reading}

\paragraph{Only the product.} In equilibrium the rotation of a scale-invariant weight depends on $\eta$ and $\lambda$ only through $\eta\lambda$ (\cref{thm:decay:equilibrium,prop:decay:adamw}); the quotient $\eta/\lambda$ sets the norm, which the function ignores. So $(\eta,\lambda)\mapsto(c\eta,\lambda/c)$ leaves the equilibrium dynamics of every scale-invariant layer unchanged, and sweeping $\lambda$ at a fixed learning rate sweeps the speed of rotation; the two settings with equal products in \cref{fig:decay:equilibrium} (middle) give the same angles. Embeddings, output layers, gains and biases see $\eta$ and $\lambda$ separately.

\paragraph{Sizes and times.} At $\eta=3\times10^{-4}$ and $\lambda=0.1$ the equilibrium angle is $\arccos(1-3\times10^{-5})=7.75\times10^{-3}$ radians per step; with $\beta_1=0.9$ it is about $1.8\times10^{-3}$ per step by \eqref{eq:decay:momentum}, while over $1000$ steps the direction turns by about $\sqrt{2\eta\lambda\cdot1000}\approx0.24$ radians either way. The norm relaxes by a factor $e$ every $1/(2\eta\lambda)\approx1.7\times10^4$ steps for AdamW, and twice as fast near equilibrium for SGD. The number of timescales $\lambda\sum_t\eta_t$ in a run varies widely: with the hyperparameters of \textcite{BrownEtAl2020} ($300$B tokens, cosine decay to $10\%$), the smallest GPT-3 model ($6\times10^{-4}$, batches of $0.5$M tokens) covers about $18$, and the largest ($6\times10^{-5}$, batches of $3.2$M tokens) about $0.3$. If its layers behaved like the idealization, the largest model never reached equilibrium.

\begin{figure}[t]
\begin{fullwidth}
\centering
\begin{tikzpicture}
\begin{groupplot}[group style={group size=3 by 1,horizontal sep=15mm},wplot,width=59mm,height=50mm]
\nextgroupplot[xlabel={$\eta\lambda t$},ylabel={$\norm{\theta_t}_2^2/q_\infty$},ymode=log,ymin=0.045,ymax=22,xmin=0,xmax=4,
  ytick={0.0625,0.25,1,4,16},yticklabels={$\frac1{16}$,$\frac14$,$1$,$4$,$16$},
  legend style={at={(0.97,0.03)},anchor=south east,font=\scriptsize}]
\addplot[theory] table[x=s,y=sgd_hi] {figures/data/13-decay-relaxation-theory.dat};
\addplot[theory2] table[x=s,y=sgn_hi] {figures/data/13-decay-relaxation-theory.dat};
\addplot[theory,forget plot] table[x=s,y=sgd_lo] {figures/data/13-decay-relaxation-theory.dat};
\addplot[theory2,forget plot] table[x=s,y=sgn_lo] {figures/data/13-decay-relaxation-theory.dat};
\addplot[simulation,forget plot] table[x=s,y=sgd_hi] {figures/data/13-decay-relaxation-sim.dat};
\addplot[simulation2,forget plot] table[x=s,y=sgn_hi] {figures/data/13-decay-relaxation-sim.dat};
\addplot[simulation,forget plot] table[x=s,y=sgd_lo] {figures/data/13-decay-relaxation-sim.dat};
\addplot[simulation2,forget plot] table[x=s,y=sgn_lo] {figures/data/13-decay-relaxation-sim.dat};
\legend{SGD,sign descent}
\nextgroupplot[xlabel={$\eta\lambda$},ylabel={angle$\,/\arccos(1-\eta\lambda)$},xmode=log,xmin=3.5e-4,xmax=2.8e-2,ymin=0,ymax=1.17,
  ytick={0,0.23,0.5,0.75,1},yticklabels={$0$,$0.23$,$0.5$,$0.75$,$1$},
  legend style={at={(0.5,0.37)},anchor=south,font=\scriptsize},legend columns=2]
\addplot[theory,domain=3.5e-4:2.8e-2,samples=2,forget plot] {1};
\addplot[theory2,domain=3.5e-4:2.8e-2,samples=2,forget plot] {0.229416};
\addplot[simulation] table[x=el,y expr=\thisrow{sgd1a}/\thisrow{exact}] {figures/data/13-decay-sweep.dat};
\addplot[simulation2,mark size=1.9pt] table[x=el,y expr=\thisrow{adamz1a}/\thisrow{exact}] {figures/data/13-decay-sweep.dat};
\addplot[only marks,mark=square,mark size=1.9pt,color=seadeep] table[x=el,y expr=\thisrow{adamz2a}/\thisrow{exact}] {figures/data/13-decay-sweep.dat};
\addplot[only marks,mark=triangle*,mark size=2.3pt,color=sea,mark options={fill=sea,draw=white,line width=0.2pt}] table[x=el,y expr=\thisrow{adamm1a}/\thisrow{exact}] {figures/data/13-decay-sweep.dat};
\addplot[only marks,mark=triangle,mark size=2.3pt,color=sea] table[x=el,y expr=\thisrow{adamm2a}/\thisrow{exact}] {figures/data/13-decay-sweep.dat};
\addplot[simulation,forget plot] table[x=el,y expr=\thisrow{sgd2a}/\thisrow{exact}] {figures/data/13-decay-sweep.dat};
\addplot[simulation,forget plot] table[x=el,y expr=\thisrow{sgd1b}/\thisrow{exact}] {figures/data/13-decay-sweep.dat};
\addplot[simulation,forget plot] table[x=el,y expr=\thisrow{sgd2b}/\thisrow{exact}] {figures/data/13-decay-sweep.dat};
\addplot[simulation2,mark size=1.9pt,forget plot] table[x=el,y expr=\thisrow{adamz1b}/\thisrow{exact}] {figures/data/13-decay-sweep.dat};
\addplot[only marks,mark=square,mark size=1.9pt,color=seadeep,forget plot] table[x=el,y expr=\thisrow{adamz2b}/\thisrow{exact}] {figures/data/13-decay-sweep.dat};
\addplot[only marks,mark=triangle*,mark size=2.3pt,color=sea,mark options={fill=sea,draw=white,line width=0.2pt},forget plot] table[x=el,y expr=\thisrow{adamm1b}/\thisrow{exact}] {figures/data/13-decay-sweep.dat};
\addplot[only marks,mark=triangle,mark size=2.3pt,color=sea,forget plot] table[x=el,y expr=\thisrow{adamm2b}/\thisrow{exact}] {figures/data/13-decay-sweep.dat};
\legend{SGD,{$\beta_1{=}0$, $W_1$},{$\beta_1{=}0$, $W_2$},{$\beta_1{=}0.9$, $W_1$},{$\beta_1{=}0.9$, $W_2$}}
\nextgroupplot[xlabel={$t/T$},ylabel={angle per step (rad)},ymode=log,ymin=1e-5,ymax=0.1,xmin=0,xmax=1,
  legend style={at={(0.03,0.03)},anchor=south west,font=\scriptsize}]
\addplot[theory] table[x=tT,y=cos_rec] {figures/data/13-decay-schedule.dat};
\addplot[theory3] table[x=tT,y=wsd_rec] {figures/data/13-decay-schedule.dat};
\addplot[guide] table[x=tT,y=cos_quasi] {figures/data/13-decay-schedule.dat};
\addplot[guide,forget plot] table[x=tT,y=wsd_quasi] {figures/data/13-decay-schedule.dat};
\addplot[simulation2,each nth point=4,forget plot] table[x=tT,y=cos_sim] {figures/data/13-decay-schedule.dat};
\addplot[simulation,each nth point=4,forget plot] table[x=tT,y=wsd_sim] {figures/data/13-decay-schedule.dat};
\legend{cosine,WSD,quasi-static}
\end{groupplot}
\end{tikzpicture}
\end{fullwidth}
\caption{The rotational equilibrium: exact predictions (curves) and runs (marks). Left: squared norm of the linear model $x\mapsto\ip{x,\hat\theta}$ in $\R^{1024}$, trained by SGD (circles) and by sign descent, which is AdamW without memory (squares), with decoupled decay at $\eta\lambda=10^{-3}$ and batch $8$ on labels independent of the inputs, so that the gradient on the sphere is noise with $\Gamma^2=1023\cdot2/8$; one run per start (seeds $11$, $12$); curves from \cref{thm:decay:equilibrium}(i) and \cref{prop:decay:adamw}. Middle: root-mean-square angle per step, over the last $4000$ of $10^4$ steps, divided by $\arccos(1-\eta\lambda)$, for the two scale-invariant matrices of a network with two hidden layers of width $32$ followed by layer normalization, trained on teacher labels with batch $16$; SGD (circles; both layers), AdamW with $\beta_1=0$ (squares) and with $\beta_1=0.9$ (triangles; filled: $W_1$, open: $W_2$), $\beta_2=0.95$; two $(\eta,\lambda)$ pairs per product, which overlap; lines at $1$ and at the heuristic value $\sqrt{(1-\beta_1)/(1+\beta_1)}$ of \eqref{eq:decay:momentum}. Right: angle per step of the model on the left ($d=2048$, sign steps) under cosine (squares) and warmup-free WSD (circles) schedules with $\eta_{\max}\lambda=10^{-3}$ over $T=2\times10^4$ steps: recursion of \cref{prop:decay:adamw}, one run (every fourth window of $200$ steps), and the quasi-static $\arccos(1-\eta_t\lambda)$ (dashed). Scripts: \texttt{13-decay-relaxation.py}, \texttt{13-decay-sweep.py}, \texttt{13-decay-schedule.py}.}
\label{fig:decay:equilibrium}
\end{figure}

\paragraph{The figure.} In \cref{fig:decay:equilibrium} (left) the norm relaxes as the recursions predict, within the run-to-run fluctuation of about $4\%$ expected from a random walk in $1024$ dimensions. From $q_\infty/16$, SGD arrives within $\eta\lambda t\approx0.25$ while sign descent needs the full linear relaxation; from $16q_\infty$ both first decay like $a^t$. In the middle, SGD's angle in a network with layer normalization lies within $1.5\%$ of $\arccos(1-\eta\lambda)$ for every product and within $0.1\%$ for $\eta\lambda\ge2\times10^{-3}$; the small excess at the smallest products is the drift term of \cref{prop:decay:average}, since the norm is still growing slowly there. AdamW without momentum matches on the first layer and falls $10$ to $15\%$ short on the second, where its preconditioned step has an outward component along $W_2$ that we measured to carry about a fifth of the decay budget. With $\beta_1=0.9$ the ratio is $0.23$ on the first layer, as \eqref{eq:decay:momentum} predicts, and about $0.2$ on the second. On the right the recursion follows the run within $2\%$ under both schedules, and the angle falls well below its quasi-static value, to $0.36$ of it at $90\%$ of the cosine schedule and $0.06$ at $99\%$, as \cref{prop:decay:schedule} requires.

\begin{figure}[t]
\begin{fullwidth}
\centering
\begin{tikzpicture}
\begin{groupplot}[group style={group size=2 by 1,horizontal sep=16mm},wplot,width=80mm,height=52mm]
\nextgroupplot[xlabel={$\eta\lambda t$},ylabel={$\lambda\abs{\theta_i}$},ymode=log,xmin=0,xmax=10,ymin=0.012,ymax=4,
  ytick={0.02,0.1,0.5,1,3},yticklabels={$0.02$,$0.1$,$0.5$,$1$,$3$},
  legend style={at={(0.98,0.45)},anchor=east,font=\scriptsize},legend columns=1]
\addplot[guide,domain=0:10,forget plot] {1};
\addplot[theory] table[x=s,y=envC] {figures/data/13-decay-fwball.dat};
\addplot[theory3] table[x=s,y=meanS] {figures/data/13-decay-fwball.dat};
\addplot[theory2] table[x=s,y=rmsN] {figures/data/13-decay-fwball.dat};
\addplot[simulation,each nth point=5,forget plot] table[x=s,y=sgnC] {figures/data/13-decay-fwball.dat};
\addplot[simulation,each nth point=5,forget plot] table[x=s,y=sgnS] {figures/data/13-decay-fwball.dat};
\addplot[simulation,each nth point=5,forget plot] table[x=s,y=sgnN] {figures/data/13-decay-fwball.dat};
\addplot[only marks,mark=square,mark size=2.0pt,color=seadeep,each nth point=5,forget plot] table[x=s,y=adamC] {figures/data/13-decay-fwball.dat};
\addplot[only marks,mark=square,mark size=2.0pt,color=seadeep,each nth point=5,forget plot] table[x=s,y=adamS] {figures/data/13-decay-fwball.dat};
\addplot[only marks,mark=square,mark size=2.0pt,color=seadeep,each nth point=5,forget plot] table[x=s,y=adamN] {figures/data/13-decay-fwball.dat};
\legend{consistent,{signal-to-noise $1$},noise}
\nextgroupplot[xlabel={step $t$},ylabel={$\Loss(\theta_t)-\Loss^\star_\lambda$},xmode=log,ymode=log,xmin=1,xmax=1e5,ymin=1e-10,ymax=30,
  legend style={at={(0.03,0.03)},anchor=south west,font=\scriptsize},legend columns=1]
\addplot[theory3] table[x=t,y=boundfw] {figures/data/13-decay-fwgap.dat};
\addplot[theory2] table[x=t,y=boundc2] {figures/data/13-decay-fwgap.dat};
\addplot[theory] table[x=t,y=boundc3] {figures/data/13-decay-fwgap.dat};
\addplot[simulation,forget plot] table[x=t,y=fw] {figures/data/13-decay-fwgap.dat};
\addplot[simulation2,forget plot] table[x=t,y=c2] {figures/data/13-decay-fwgap.dat};
\addplot[only marks,mark=triangle*,mark size=2.0pt,color=sea,mark options={fill=sea,draw=white,line width=0.2pt},forget plot] table[x=t,y=c3] {figures/data/13-decay-fwgap.dat};
\legend{$\gamma_t=2/(t+2)$,$\gamma=10^{-2}$,$\gamma=10^{-3}$}
\end{groupplot}
\end{tikzpicture}
\end{fullwidth}
\caption{AdamW and the Frank--Wolfe ball $\norm{\theta}_\infty\le1/\lambda$. Left: sign descent (circles) and AdamW with $(\beta_1,\beta_2)=(0.9,0.95)$ (open squares), $\eta=0.01$, $\lambda=0.1$, on a linear loss with Gaussian gradient noise, every coordinate started at $3/\lambda$ (seed $31$): $20$ coordinates with a consistent gradient (largest $\lambda\abs{\theta_i}$; curve: the equality case of \cref{thm:decay:frankwolfe}(ii)), $200$ with signal-to-noise ratio $1$ (mean; curve: \cref{prop:decay:regimes}(i), limit $\erf(1/\sqrt2)=0.683$) and $2000$ with pure noise (root mean square; curve: \cref{prop:decay:regimes}(ii), limit $0.0224$). Right: the constrained gap of sign descent with decay on a convex quadratic in $\R^{10}$ with $\lambda=1$ and five active constraints at the optimum (seed $7$; $L=6.35$ by enumeration of sign vectors); marks: runs with the step rule of the curve of the same color; curves: the guarantees of \cref{thm:decay:frankwolfe}(iii). Script: \texttt{13-decay-frankwolfe.py}.}
\label{fig:decay:frankwolfe}
\end{figure}

\paragraph{Where the weights sit.} In \cref{fig:decay:frankwolfe} (left), coordinates with a consistent gradient are pulled to the face of the ball exactly along $1+(1-\eta\lambda)^t(\lambda\theta_0-1)$, so the attraction rate of \cref{thm:decay:frankwolfe}(ii) is attained; coordinates with signal-to-noise ratio one settle at $0.683/\lambda$ under sign descent and $0.705/\lambda$ under AdamW, whose averaged step is slightly larger; noise coordinates settle at $0.0224/\lambda$ under both. On the right the guarantees hold with room to spare: under $2/(t+2)$ the gap falls like a few times $1/t^2$, and at a constant step it fluctuates around $0.4\gamma^2$, far below $2L\gamma/\lambda^2$.

\paragraph{Where the agreement must fail.} The equilibrium predictions are exact for orthogonal steps, so they fail where steps are not orthogonal (momentum and preconditioning, the middle panel), where the gradient on the sphere changes faster than the norm can follow (early training, and the feedback discussed next), and where training is shorter than $1/(\eta\lambda)$. The Frank--Wolfe rates are worst cases over convex $L$-smooth losses and say nothing about a network; what transfers is the ball.

\section{What the theory misses}\label{sec:decay:misses}

\paragraph{What was idealized away.} \emph{Exact scale invariance:} embeddings, output layers, gains and biases are not scale invariant, and a matrix whose output enters the residual stream before a normalization is so only approximately; for these parameters $\eta$ and $\lambda$ act separately, and the decay does regularize. \emph{A fixed gradient size on the sphere:} in a network $\Gamma_t$ falls with the loss and depends on the effective step itself, since a larger effective step keeps the direction farther from the minimizers on the sphere. In a teacher--student version of \cref{fig:decay:equilibrium} (left), the student's wandering raised $\Gamma^2$ by $17\%$ above its noise value and the norm settled $8\%$ above the prediction made with the noise value; \cref{thm:decay:equilibrium}(iv) brackets such effects but does not model the feedback. \emph{Orthogonal steps:} preconditioners and momentum add to \cref{prop:decay:average} the radial term $-2\eta(1-\eta\lambda)\sum_t\ip{u_t,\theta_t}$, which in our network carried a fifth of the decay budget without momentum and most of it with momentum. \emph{Convexity:} the rates of \cref{thm:decay:frankwolfe} need a convex loss with a useful smoothness constant in the chosen norm; the ball statements need neither.

\paragraph{Where it breaks.} (1) The equilibrium need not be stable. \Textcite{LobachevaEtAl2021} found periodic cycles with batch normalization and weight decay, and \textcite{AminChangKhanna2026} (unrefereed, 2026) argue that a constant learning rate with weight decay cannot hold a stable interior equilibrium in such networks. By \cref{lem:decay:scaleinv}(iii), a shrinking norm raises the curvature that the step meets, so the instability of \cref{ch:eos} can return in this form; with $\Gamma$ fixed, our theorem cannot see it. (2) In vision networks weight decay does more than set a rotation speed: \textcite{ZhangEtAl2019} attribute part of its effect on test error to other mechanisms, and nothing here predicts test error. (3) The constrained solutions of the Frank--Wolfe view are not where pretraining ends: at $\lambda=0.1$ the ball is far larger than the weights, and the limit theorems of \textcite{XieLi2024} and \textcite{ChenEtAl2024Lion} assume convergence, which a finite schedule does not deliver; the constraint matters for consistent coordinates and for the implicit bias of long training (\cref{ch:bias}). (4) At scale, the effective-learning-rate view is supported but not settled. \Textcite{LiuEtAl2026ELR} (unrefereed, 2026) report that loss curves of language models with different learning rates and norms collapse, to within a few thousandths and below the variation between seeds, when indexed by the effective learning rate, and that weight decay and Hyperball act mainly through the effective-learning-rate schedules they induce. \Textcite{WenEtAl2026Hyperball} (unrefereed, 2026) report that fixing the norms with Hyperball speeds up Muon by $20$ to $30\%$ in token-equivalent terms over weight-decay baselines up to $1.2$B parameters and reduces the drift of the optimal learning rate across depths, which suggests that the equilibrium is not the best way to set the effective learning rate; \textcite{XiaoEtAl2026} (unrefereed, 2026) trace the difference between Muon with Hyperball and Muon with decay to the time course of the effective step size rather than to a better update direction, and \textcite{YuanEtAl2026} (unrefereed, 2026) show that for scale-invariant networks the two are dynamically equivalent, as in \cref{prop:decay:hyperball}. (5) Equilibrium takes time: for the largest runs the transient, not the equilibrium, may be the regime of interest (\cref{sec:decay:reading}).

\paragraph{Repairs and open problems.} Fixing the norm removes the transient and makes the angle a hyperparameter: Hyperball, and the rotational variants of AdamW, SGD and Lion of \textcite{KossonMessmerJaggi2024}, which control the rotation explicitly and reduce the need for warmup (\cref{ch:warmup}). \Cref{prop:decay:hyperball} says what such a method must do to imitate weight decay, and \cref{prop:decay:schedule} what it gains: an angle that follows the schedule without lag. Scaling rules for $\lambda$ \parencite{LoshchilovHutter2019,WangAitchison2024} rest on hypotheses about the right number of timescales that no theory derives. Open: a model in which $\Gamma$ depends on the effective step and the stability of the equilibrium, or the period of the cycles, can be computed (\cref{ex:decay:feedback}); a theory of decay for parameters that are not scale invariant; and a joint theory of decay and schedule in the river valley of \cref{ch:schedules} (\cref{ex:decay:river}).

\section{Back in the wild}\label{sec:decay:back}

\fieldnotes{AdamW with decoupled decay $\lambda=0.1$ at $\eta\sim10^{-4}$; decay for Muon at scale; rescaling $\lambda$ with training length and width.}
{(a) Steps of fixed length in a norm, decoupled decay, convex losses for the rates. (b) A scale-invariant loss, steps orthogonal to the weights, constant gradient size on the unit sphere.}
{Decay is Frank--Wolfe on the ball of radius $1/\lambda$: invariance, attraction at rate $(1-\eta\lambda)^t$, rate $O(1/t)$, an $O(\eta/\lambda)$ neighborhood. The norm of a scale-invariant weight converges from every start, and the angle per step tends to $\arccos(1-\eta\lambda)$; it has that mean whenever the norm is stationary. The rotation lags the schedule; decay is a Hyperball schedule; AdamW is an EMA of its updates.}
{Only $\eta\lambda$ matters in equilibrium: $7.7\times10^{-3}$ rad per step at $\eta\lambda=3\times10^{-5}$, after about $1/(2\eta\lambda)$ steps; momentum divides the angle per step by $\sqrt{(1+\beta_1)/(1-\beta_1)}$ (derived); a decaying schedule cools faster than $\sqrt{\eta_t\lambda}$.}
{Layers that are not scale invariant; gradients that respond to the effective step; radial components of preconditioned and momentum steps; generalization; periodic instabilities.}
{Treat $\eta\lambda$ as the hyperparameter of scale-invariant layers and keep it when $\eta$ is rescaled (proved in the idealization); keep $\eta\lambda T$ fixed when the run is lengthened (suggested, with evidence); expect the effective learning rate to lag the schedule (proved in the model).}

The theory tells a practitioner what $\lambda$ does, not what it should be. For the layers that feed a normalization, the decay is a speed: once the norm has settled, the weights turn by about $\sqrt{2\eta\lambda}$ radians per step whatever the gradients, so a sweep over $\lambda$ is a sweep over that speed, and a change of learning rate that keeps $\eta\lambda$ fixed changes nothing after the transient. The ball of radius $1/\lambda$ is real but, at $\lambda=0.1$, far from the weights of a language model; it shapes where long training ends (\cref{ch:bias}) more than how pretraining proceeds. Which product $\eta\lambda$ is best, how many timescales a run should cover, and whether an equilibrium beats a fixed norm are empirical questions, and the 2026 evidence on them is unrefereed.

The chapter closes several threads. The Frank--Wolfe identity is the normalized steepest descent of \cref{ch:steepest}, in the norms of \cref{ch:sign,ch:spectral}, made convergent by its schedule, the thread of \cref{ch:momentum,ch:schedules}. The equilibrium angle is an effective learning rate: the per-layer rates of $\mu$P (\cref{ch:modular}) and the angular updates of warmup (\cref{ch:warmup}) seen from the sphere. And by \cref{lem:decay:scaleinv}(iii), the stability threshold of \cref{ch:eos} for a scale-invariant layer is a threshold on $\eta/\norm{\theta}_2^2$, which the decay controls.

\begin{exercises}
\item \consolidation (Frank--Wolfe on an interval.) Let $f(\theta)=\frac h2\theta^2$ and $\theta_{t+1}=(1-\gamma_t)\theta_t-\gamma_t\lambda^{-1}\sign(\theta_t)$. (a) For a constant $\gamma\in(0,1)$, show that the two-cycle $\pm\gamma/(\lambda(2-\gamma))$ attracts every start, so that $f$ settles at $\frac h2\gamma^2/(\lambda^2(2-\gamma)^2)$, about $\gamma/16$ times the guarantee $2h\gamma/\lambda^2$. (b) For $\gamma_t=2/(t+2)$ and $\theta_0\ne0$, show that $\abs{\theta_t}=1/(\lambda(2\lfloor t/2\rfloor+1))$ for $t\ge1$.\label{ex:decay:cycle}

\item \consolidation (The Frank--Wolfe gap.) For convex differentiable $f$ and $\theta\in\decayball$, show that $\mathcal G(\theta)=\ip{\nabla f(\theta),\theta}+\norm{\nabla f(\theta)}_*/\lambda$ is at least $f(\theta)-\min_{\decayball}f$, and vanishes exactly at the minimizers over $\decayball$. For the $\ell_\infty$ ball, deduce that at a minimizer every coordinate has zero partial derivative or sits at $-\sign(\partial_if)/\lambda$.\label{ex:decay:gap}

\item \consolidation (Schedules for Frank--Wolfe.) Let $h_{t+1}\le(1-\gamma_t)h_t+C\gamma_t^2$ with $\gamma_t\in[0,1]$, $\gamma_t\to0$ and $\sum_t\gamma_t=\infty$. Show that $\limsup_th_t\le0$. Compute the bound on $h_T$ for the linear decay $\gamma_t=\gamma_0(1-t/T)$ and compare it with the constant step $\gamma_0$.\label{ex:decay:schedules}

\item \consolidation (The scalar map.) For $\varphi(x)=ax+(1-a)/x$: (a) verify \eqref{eq:decay:factor}; (b) show that for $a<\frac12$ the iterates approach $1$ from alternating sides; (c) show that for $a=0$ every point lies on a two-cycle; (d) for $a=(1-\eta\lambda)^2$ and $x_0=e^{10}$, show that about $10/(2\eta\lambda)$ steps are needed to reach $x_t\le e$.\label{ex:decay:map}

\item \consolidation (Angles.) Show that $\sin\omega_\infty=\sqrt{\eta\lambda(2-\eta\lambda)}$ and $2\sin(\omega_\infty/2)=\sqrt{2\eta\lambda}$. At $\eta=0.05$ and $\lambda=0.01$, evaluate $\omega_\infty$, $\sin\omega_\infty$ and $\sqrt{2\eta\lambda}$ to six digits, and explain why a simulation of the angle gives $0.031624$.\label{ex:decay:angles}

\item \variation (Adam with an $L_2$ penalty.) In Adam with $\beta_1=0$, replace decoupled decay by a penalty: $u_t=(g_t+\lambda\theta_t)/\sqrt{\hat v_t}$, $\theta_{t+1}=\theta_t-\eta u_t$, on a scale-invariant layer whose gradient on the unit sphere is noise of variance $\sigma^2$ per coordinate. Assuming $\hat v_t\approx\sigma^2/\norm{\theta_t}_2^2$, show that the decay per step is $\eta\lambda\norm{\theta_t}_2/\sigma$ and that the norm settles at $\norm{\theta}_2^3\approx\eta\sigma d/(2\lambda)$, so that the equilibrium angle depends on the size of the gradient, unlike AdamW's \parencite{KossonMessmerJaggi2024}. (A simulation with $\beta_2=0.999$ agrees to $0.3\%$.)\label{ex:decay:l2}

\item \variation (Momentum, exactly.) Let $\theta_{t+1}=c\theta_t-\eta m_t$ and $m_t=\beta m_{t-1}+(1-\beta)\xi_t$ with $c=1-\eta\lambda$ and independent $\xi_t$ of mean zero and covariance $I_d$. (a) Using \cref{lem:toolkit:lyapunov} for the pair $(\theta_{t,i},m_{t-1,i})$, show that the stationary $\E\norm{\theta}_2^2$ is $\eta^2d\,(1-\beta)(1+c\beta)/((1+\beta)(1-c\beta)(1-c^2))$, which is $\eta^2d/(1-c^2)$ times $1+O(\eta\lambda/(1-\beta))$. (b) Show that $\E\norm{m}_2^2=d(1-\beta)/(1+\beta)$ and deduce \eqref{eq:decay:momentum}. (c) For PyTorch's SGD with momentum, $m_t=\beta m_{t-1}+g_t+\lambda\theta_t$ and $\theta_{t+1}=\theta_t-\eta m_t$, show in the same model that the angle per step is about $\sqrt{2\eta\lambda/(1+\beta)}$.\label{ex:decay:momentum}

\item \variation (MuonW.) A scale-invariant $d_{\mathrm{out}}\times d_{\mathrm{in}}$ matrix is trained by $W_{t+1}=(1-\eta\lambda)W_t-\eta s\polar(G_t)$ with full-rank gradients $G_t$. (a) Show with a $2\times2$ example that $\ip{\polar(G_t),W_t}$ need not vanish although $\ip{G_t,W_t}=0$. (b) Assuming it averages to zero, compute the equilibrium $\norm{W}_F$, and compare the spectral norm of a Gaussian matrix with that Frobenius norm with the radius $s/\lambda$ of the Frank--Wolfe ball, for $s=0.2\sqrt{\max\{d_{\mathrm{out}},d_{\mathrm{in}}\}}$, $\eta=10^{-3}$, $\lambda=0.1$ and $d_{\mathrm{in}}=d_{\mathrm{out}}=1024$.\label{ex:decay:muon}

\item \variation (The lag under linear decay.) In \cref{prop:decay:adamw}(i), let $\eta_t=\eta_0(1-t/T)$ with $\eta_0\lambda T\gg1$, started at $q_\infty(\eta_0)$. Treating $t$ as continuous, show that the norm tracks $q_\infty(\eta_t)$ while $\eta_t\lambda(T-t)\gg1$ and freezes once this is of order one, after which the angle falls like $\eta_t$. Show that the transition happens when the remaining fraction of the run is about $(\eta_0\lambda T)^{-1/2}$.\label{ex:decay:lag}

\item \simulation (Universality.) Run \texttt{13-decay-sweep.py} with $\beta_1=0.5$, with $\beta_2=0.999$, and with width $64$. Check that SGD's angle stays at $\arccos(1-\eta\lambda)$ and that AdamW's with $\beta_1=0.5$ is close to $1/\sqrt3$ of it on the first layer, as \eqref{eq:decay:momentum} predicts; measure the radial component of AdamW's step on each layer and explain why it can differ between layers.\label{ex:decay:sim1}

\item \simulation (Where in the ball.) In \texttt{13-decay-frankwolfe.py}, vary the signal-to-noise ratio of the middle group over $\{0.25,0.5,1,2,4\}$; check that sign descent settles at $\erf(\mathrm{snr}/\sqrt2)/\lambda$, measure where AdamW settles, and explain why the two differ most at intermediate ratios. With $\eta=3\times10^{-4}$ and $\lambda=0.1$, check that the noise coordinates settle at $0.39\%$ of the radius.\label{ex:decay:sim2}

\item \research (Feedback and stability.) Let the gradient size on the sphere depend on the effective step, $\Gamma^2=\Gamma_0^2+\kappa\,\decayeff$, as for a student that wanders around a teacher. Find the equilibrium of \cref{thm:decay:equilibrium} and its stability. Is there a natural model, for instance one in which $\decayeff S(\hat\theta)$ is capped at $2$ as in \cref{ch:eos}, where the equilibrium loses stability and the cycles of \textcite{LobachevaEtAl2021} appear? As far as we know this is open.\label{ex:decay:feedback}

\item \research (Weight decay in the river valley.) Combine the river-valley model of \cref{ch:schedules} with a scale-invariant layer whose direction carries the river and the hills. For a given schedule, which $\lambda$ minimizes the final loss, and how does the answer depend on the length of training? Compare with the timescale rule of \cref{prop:decay:ema}.\label{ex:decay:river}
\end{exercises}

\notesandrefs
Decoupled weight decay is due to \textcite{LoshchilovHutter2019}, who showed that it differs from an $L_2$ penalty for Adam \parencite{KingmaBa2015}. That normalization turns a penalty on the weights into a control of the effective learning rate was observed by \textcite{VanLaarhoven2017}, and \textcite{ZhangEtAl2019} placed it among three mechanisms of weight decay. \Textcite{WanEtAl2021} derived the equilibrium angle for SGD with momentum and proved, under their assumptions, convergence to the equilibrium at a linear rate; \textcite{KossonMessmerJaggi2024} measured it per neuron for several optimizers, and \textcite{LobachevaEtAl2021} found the periodic alternative. \Cref{thm:decay:equilibrium} is the deterministic core of these analyses, with convergence from every initialization and a comparison bound; we know no earlier statement of the identity of \cref{prop:decay:average} in this form.

The Frank--Wolfe method was introduced by \textcite{FrankWolfe1956} for quadratic programs; \textcite{Jaggi2013} gave the analysis with the open-loop step $2/(t+2)$ that made it popular in machine learning, and \textcite{Hazan2016} treats it among projection-free methods. The identification of decoupled decay with Frank--Wolfe on a norm ball runs through \textcite{XieLi2024} for AdamW, \textcite{ChenEtAl2024Lion} for Lion, \textcite{ChenLiLiu2025} for Muon and \textcite{PethickEtAl2025} for general oracles; \textcite{BernsteinNewhouse2024} cast these optimizers as steepest descent under norms, and \textcite{Newhouse2025} connects duality, weight decay and norm constraints in a thesis. The moving-average reading of AdamW and its scaling rules are those of \textcite{WangAitchison2024}. The 2026 literature on fixed norms and effective learning rates is unrefereed \parencite{WenEtAl2026Hyperball,LiuEtAl2026ELR,XiaoEtAl2026,YuanEtAl2026,AminChangKhanna2026}. The figures were computed for this chapter; the network of \cref{fig:decay:equilibrium} is small, and its agreement with the theory checks the identities, not the behavior of large models.

\paragraph{Translating the literature.} \Textcite{LoshchilovHutter2019} write AdamW as $\theta_t=\theta_{t-1}-\eta_t\bigl(\alpha\hat m_t/(\sqrt{\hat v_t}+\epsilon)+\lambda\theta_{t-1}\bigr)$, with $\alpha$ the learning rate and $\eta_t$ a schedule multiplier: their $\alpha\eta_t$ is our $\eta_t$ and their $\lambda$ is our $\alpha\lambda$. PyTorch's AdamW decays by \texttt{lr}$\cdot$\texttt{weight\_decay} per step and is our $(\eta,\lambda)$ \parencite{Schaipp2024}. \Textcite{WangAitchison2024} write $\tau_{\mathrm{iter}}=1/(\eta\lambda)$ for our $\tau$, and $\tau_{\mathrm{epoch}}=\tau_{\mathrm{iter}}/M$ with $M$ iterations per epoch. \Textcite{KossonMessmerJaggi2024} follow the weight vector of one neuron, a row of our $W$; their angular update, like that of \textcite{WanEtAl2021}, is our $\omega_t$. Hyperball's weights have norm $R$, and its step of norm $\eta_tR$ before renormalization makes its $\eta_t$ our $\alpha_t$ of \cref{prop:decay:hyperball}. \Textcite{LiuEtAl2026ELR} define the effective learning rate as the learning rate divided by the parameter norm, proportional to our $\tan\omega_t$ for steps of fixed size; for SGD our effective step size divides by the squared norm. In the Frank--Wolfe literature $\gamma_t$ is the step size and $\cK$ the constraint set, here $\eta_t\lambda$ and $\decayball$.

\printchapterbib
\end{refsection}
